\section{Conclusion}

We set out to test whether the scaling principles established for general-domain agent systems
hold in medicine, in the one regime those systems cannot reach: tasks with no tools. Three results
follow.

The capability ceiling transfers, but as one edge of a two-sided window. Collaboration pays when a
single physician model scores roughly 25--50\% and not otherwise, and this is a sharper and more
useful statement than a ceiling because it identifies both failure modes: on easy cases there is
nothing to gain, and on the hardest cases a panel of unreliable specialists reinforces its own
errors rather than correcting them. The general-domain data are consistent with a monotone ceiling
only because their benchmarks never descend into the difficult regime.

Coordination structure dominates panel size. Panel size does not predict collaboration gain at
all, gains saturate at three agents, and the only significant improvements come from centralized
coordination in which an attending physician audits specialist opinions before committing. Where
coordination has no verification channel, error amplification --- the single strongest term in our
scaling model --- goes uncontrolled and the panel discards correct answers it already held.

The economics do not favour consultation outside the window. A panel costs between $2\times$ and
$5\times$ a single call per correct answer, and outside the window a budget-matched
self-consistency control matches or beats every architecture we tested. For medical deployments
this converts an architectural question into a triage question: measure the single-model baseline
on the case mix first, and consult only where the window says consultation will pay.
